%! Introduction to Functions and Change — Unit 1, Lesson 1.0 — Exit Ticket (Answer Key)
\documentclass[10pt]{article}
\usepackage{precalculus-article}
\usepackage{precalculus-key}   % pulls in -boxes; do NOT also load -boxes
\newcommand{\TallMath}[1]{$\displaystyle #1\rule[-1.4em]{0pt}{3.2em}$}

\begin{document}

\pageheader{Unit 1, Lesson 1.0}{Exit Ticket --- Answer Key}

\vspace{0.12in}

\begin{enumerate}[itemsep=12pt]

\item An oven is preheating, and we track its temperature as the minutes pass.

\vspace{0.06in}
Input quantity: \ans{time (minutes)} \quad Output quantity: \ans{oven temperature}

\item Consider the rule: each car $\to$ its license plate number.
Is this rule a function? Circle one and explain in one sentence.

\vspace{0.06in}
\ans{Yes} \quad / \quad \textbf{No}

\vspace{0.02in}
\ansline{Each car has exactly one license plate number, so each input has exactly one output.}

\item The table shows a plant's height over several weeks.

\smallskip
\begin{tabular}{|c|c|c|c|c|}
\hline
Week & 0 & 2 & 4 & 6 \\
\hline
Height (in) & 3 & 3 & 5 & 8 \\
\hline
\end{tabular}
\smallskip

\begin{enumerate}[label=(\alph*), itemsep=4pt]
  \item The height is \textbf{constant} from week \ans{0} to week \ans{2}.
  \item The height is \textbf{increasing} from week \ans{2} to week \ans{6}.
\end{enumerate}

\end{enumerate}

\end{document}
